\chapter{Device drivers, the registry, and how an abort reaches a call}\label{ch:devices}

A Device is anything a Test Schedule opens: a TCP socket, a UDP peer, a serial port, a VISA resource.
Each Device Type has one driver, and the driver contract is small on purpose --- a new Device Type is a
file and one line, and chapter~\ref{ch:device-type} walks through it with a driver that the test suite
actually runs.

\section{The contract}

\texttt{ate/\allowbreak{}devices/\allowbreak{}registry.py} defines \texttt{Driver}, an abstract base class with three class
attributes and two methods. Only \texttt{open} is abstract. Teardown is not the driver's job: the handle
\texttt{open} returns has its own \texttt{close}, and that is what the registry calls. \texttt{Driver.close}
is declared with a no-op default, but nothing reaches it --- none of the four shipped drivers is
teared down through the driver, and \texttt{VisaDriver.close} is the only override, dead code kept so the
contract reads as written.

\begin{lstlisting}[language=ATEPython]
class Driver(ABC):
    name: ClassVar[str] = "driver"                       # the Device Type name, e.g. "TCP"
    accepts_receive_buffer_bytes: ClassVar[bool] = False # does open() take the buffer cap?
    fields: ClassVar[dict[str, tuple[tuple[type, ...], bool]]] = {}  # field -> (types, required)

    @abstractmethod
    def open(self, params: dict[str, Any]) -> Any: ...

    def close(self, device: Any) -> None:
        """Release a device. The registry closes the handle, never through this method."""
\end{lstlisting}

\texttt{validate\_params} checks every declared field generically: an unknown field, a missing required
one, or a value of the wrong type raises a \texttt{ValueError} naming the field. Booleans are explicitly
not numbers --- \texttt{isinstance(True, int)} is true in Python, so a driver that accepts
\texttt{(int, float)} must reject a boolean, and \texttt{\_matches} does that.

The registry never knows how a driver connects. It only knows that \texttt{open} produced a handle, and
it tracks that handle until \texttt{close\_device} is called on it, so a Device left open at the end of a
run can be closed and warned about (packet 16a). \texttt{close\_device} is the whole of the teardown
contract: it calls the handle's own \texttt{close}, through \texttt{run\_abortable} when the close is one
the abort is meant to end, and \texttt{ATE.close} and the run's end-of-run cleanup both go through it.
\texttt{\_StreamDevice.close} (\texttt{stream.py}) stops the reader thread, closes the transport and joins
the reader; \texttt{\_VisaDevice.close} (\texttt{visa.py}) closes the resource. Both are safe to call
twice, so a Device closed by the schedule and then closed again by the run's cleanup is not an error.

The four drivers that ship, and the Device Types they serve, are \texttt{TcpDriver} (TCP),
\texttt{SerialDriver} (SERIAL), \texttt{UdpDriver} (UDP) and \texttt{VisaDriver} (VISA).
\texttt{tests/\allowbreak{}test\_\allowbreak{}manuals.py} derives that list by importing every module in \texttt{ate.devices} and
collecting the classes that subclass \texttt{Driver}, so a driver added to the code and not named in the
manual fails the suite.

TCP, SERIAL and UDP share \texttt{devices/stream.py}: a background reader thread, a receive buffer
capped by \texttt{devices.receiveBufferBytes} (default 16777216 bytes), an overflow that drops the oldest
bytes and reports \texttt{overflow = true} with one WARNING on the next read
(\texttt{receive buffer overflow - dropped 4096 oldest byte(s)}), and the failure and warning logic.
UDP is that same machinery with one change: \texttt{\_UdpDevice} subclasses \texttt{\_StreamDevice} and
replaces \texttt{\_reader\_loop} (\texttt{udp.py}, "The shared drain loop, except that an empty datagram
is data"), because the base loop skips a falsy chunk and a zero-length datagram is a real message. VISA
is not on the stream at all: \texttt{\_VisaDevice} in \texttt{visa.py} is its own class, reads on demand
and has no reader thread, so \texttt{accepts\_\allowbreak{}receive\_\allowbreak{}buffer\_\allowbreak{}bytes} is \texttt{False} for it --- there
is no buffer to cap.

\section{Tracing}

When \texttt{deviceTrace.enabled} is on, the registry wraps every handle it opens in
\texttt{TracedDevice} and logs each open, read, write and close as one \texttt{DEBUG} line with source
\texttt{device}, through the host's log sink --- so a trace line is owned by the step or Action that made
the call, exactly like a schedule log line. Payloads are escaped and truncated at
\texttt{deviceTrace.maxBytes}, default \texttt{DEFAULT\_\allowbreak{}TRACE\_\allowbreak{}MAX\_\allowbreak{}BYTES = 256}. A call the abort left
behind gets its own line, \texttt{open <device> aborted (ms)} or \texttt{close <device> aborted (ms)}
(packet 32b review round: before the fix the trace list was empty).

\section{Abort, and what an abort cannot stop}

Story 83 names two \emph{allowed} exceptions to the ten-second bound, and the suite measures both
rather than trying to fix them: a command a step runs with \texttt{os.execute}, which the abort
cannot end, and a busy loop inside a coroutine, where the count hook never ticks
(\texttt{tests/test\_\allowbreak{}manuals.py}: \texttt{test\_\allowbreak{}a\_\allowbreak{}command\_\allowbreak{}in\_\allowbreak{}a\_\allowbreak{}step\_\allowbreak{}holds\_\allowbreak{}the\_\allowbreak{}abort\_\allowbreak{}until\_\allowbreak{}it\_\allowbreak{}finishes},
\texttt{test\_\allowbreak{}a\_\allowbreak{}coroutine\_\allowbreak{}loop\_\allowbreak{}runs\_\allowbreak{}to\_\allowbreak{}completion\_\allowbreak{}after\_\allowbreak{}an\_\allowbreak{}abort}).
Neither is forbidden in a schedule. \texttt{os.exit} is not one of them: it ends the process with the
run row left open, and the next launch closes it as \texttt{INTERRUPTED}/\texttt{FAIL}.

\texttt{ate/devices/abort.py} is packet 32b's answer to ``the abort does not stop the worker in 10 s''.
There are two kinds of Device wait:

\begin{itemize}
  \item \textbf{A wait ATE owns} --- a buffer read, a \texttt{select}, a pyvisa-py SOCKET read or write.
    It runs in slices of \texttt{POLL\_SECONDS} ($0.05$ s) and checks the flag between them, so
    it ends in about 50 ms. \texttt{wait\_slice} caps the next slice by the call's own remaining timeout,
    and \texttt{check\_abort} raises \texttt{DeviceAborted} which the ATE wrappers turn into Lua's
    \texttt{aborted} error.
  \item \textbf{A wait only the OS or a C library can end} --- a connect, a DNS lookup, \texttt{viRead} on
    a real VISA stack, a pyserial write, a close on a stuck tty. These run on their own daemon thread
    through \texttt{run\_abortable}. The worker watches the flag and raises, leaving the call behind; it
    finishes when its own timeout says so, and its result is discarded rather than recorded.
\end{itemize}

Two details that were found by test and are easy to break again. \texttt{\_Call.\_run} catches
\texttt{BaseException}, not \texttt{Exception}, because a helper thread killed by \texttt{KeyboardInterrupt}
or \texttt{SystemExit} would otherwise leave no outcome and the worker would wait forever
(\texttt{test\_\allowbreak{}a\_\allowbreak{}helper\_\allowbreak{}thread\_\allowbreak{}dying\_\allowbreak{}on\_\allowbreak{}a\_\allowbreak{}base\_\allowbreak{}exception\_\allowbreak{}reaches\_\allowbreak{}the\_\allowbreak{}worker}). And a helper thread
left behind is joined in slices, so an aborted run's leftover closes are bounded:
\texttt{\_\allowbreak{}ABORT\_\allowbreak{}CLOSE\_\allowbreak{}SECONDS = 3.0} in \texttt{run\_engine.py} is the total for all closes at the end of
an aborted run, each with a WARNING naming the Device.

The consequence for a Station, stated in the admin manual as well: an aborted run can leave a Device
close running for seconds, a Post-test Action on the same Device waits for it, and a new run opening the
same instrument meanwhile may find it busy. That is not a bug to fix; it is the boundary of what an abort
can do.

\section{Which threads a Device runs on}

The background readers are \texttt{threading.Thread} objects, not \texttt{QThread}s. That matters twice:
their exceptions go to \texttt{threading.excepthook}, not \texttt{sys.excepthook}, so the
exception-log hook does not see them --- they catch their own --- and the architecture guard does not
scan them because they are not under \texttt{src/ate/engine} or \texttt{src/ate/run\_session}.

\cando{write a driver that the registry will accept, and explain what happens to an open Device when a run
is aborted.}
